\section{传播子}
\subsection{定义、演化与合成}
在坐标表象中定义传播子
\begin{equation}
\boxed{K\left(x,t;x_0,t_0\right)
=\Braket{x|\hat U\left(t,t_0\right)|x_0}.}
\label{eq:propagator-definition}
\end{equation}
插入坐标完备性后，任意初始波函数按
\begin{equation}
\psi\left(x,t\right)=\int_{-\infty}^{\infty}
K\left(x,t;x_0,t_0\right)\psi\left(x_0,t_0\right)\,\mathrm{d}x_0
\label{eq:wavefunction-propagation}
\end{equation}
演化。由演化算符方程，传播子满足
\begin{equation}
\mathrm{i}\hbar\frac{\partial}{\partial t}K\left(x,t;x_0,t_0\right)
=\hat H\left(x,-\mathrm{i}\hbar\frac{\partial}{\partial x},t\right)
K\left(x,t;x_0,t_0\right),
\end{equation}
初值条件为
\begin{equation}
K\left(x,t_0;x_0,t_0\right)=\delta\left(x-x_0\right).
\end{equation}
演化算符的合成律对应传播子的卷积：对 $t_0<t_1<t_2$，
\begin{equation}
\boxed{K\left(x_2,t_2;x_0,t_0\right)=
\int_{-\infty}^{\infty}K\left(x_2,t_2;x_1,t_1\right)
K\left(x_1,t_1;x_0,t_0\right)\,\mathrm{d}x_1.}
\label{eq:propagator-composition}
\end{equation}

\subsection{定态谱表示}
若 Hamilton 算符不含时并具有一组离散正交完备本征态，插入能量本征态单位分解可得
\begin{equation}
\boxed{K\left(x,t;x_0,t_0\right)=
\sum_n\mathrm{e}^{-\mathrm{i}E_n\left(t-t_0\right)/\hbar}
\psi_n\left(x\right)\psi_n^*\left(x_0\right).}
\label{eq:propagator-spectrum}
\end{equation}
若存在连续谱，还要加入相应的谱积分。沿对角线积分得到演化算符的迹：
\begin{equation}
\int_{-\infty}^{\infty}K\left(x,t;x,t_0\right)\,\mathrm{d}x
=\Tr\,\hat U\left(t,t_0\right)
=\sum_n\mathrm{e}^{-\mathrm{i}E_n\left(t-t_0\right)/\hbar}.
\end{equation}

\subsection{自由粒子传播子}
对 $\hat H=\hat p^2/\left(2m\right)$，插入动量完备性，并令 $T=t-t_0>0$：
\begin{align}
K_0\left(x,t;x_0,t_0\right)
&=\int_{-\infty}^{\infty}\frac{\mathrm{d}p}{2\pi\hbar}
\exp\left[\frac{\mathrm{i}}{\hbar}p\left(x-x_0\right)
-\frac{\mathrm{i}}{\hbar}\frac{p^2}{2m}T\right]\notag\\
&=\boxed{\sqrt{\frac{m}{2\pi\mathrm{i}\hbar T}}
\exp\left[\frac{\mathrm{i}m\left(x-x_0\right)^2}{2\hbar T}\right].}
\label{eq:free-propagator}
\end{align}
平方根的分支由 $T\to0^+$ 时传播子趋于 $\delta\left(x-x_0\right)$ 的条件固定。

\subsection{谐振子传播子}
对一维谐振子，利用式\eqref{eq:propagator-spectrum}与 Hermite 多项式的 Mehler 求和公式可得
\begin{equation}
\boxed{
K_{\mathrm{ho}}\left(x,t;x_0,t_0\right)
=\sqrt{\frac{m\omega}{2\pi\mathrm{i}\hbar\sin\left(\omega T\right)}}
\exp\left\{\frac{\mathrm{i}m\omega}{2\hbar\sin\left(\omega T\right)}
\left[\left(x^2+x_0^2\right)\cos\left(\omega T\right)-2xx_0\right]\right\},
}
\label{eq:oscillator-propagator}
\end{equation}
其中 $T=t-t_0$，并暂避开 $\sin\left(\omega T\right)=0$ 的焦散时刻；跨过焦散时需要连续跟踪平方根的相位。

当 $\omega\to0$ 时，使用 $\sin\left(\omega T\right)\sim\omega T$ 与
$\cos\left(\omega T\right)\sim1$，式\eqref{eq:oscillator-propagator}回到自由粒子传播子。
另一方面，对角线积分为
\begin{equation}
\int_{-\infty}^{\infty}K_{\mathrm{ho}}\left(x,t;x,t_0\right)\,\mathrm{d}x
=\sum_{n=0}^{\infty}\mathrm{e}^{-\mathrm{i}\omega\left(n+1/2\right)T},
\end{equation}
从中可读出 $E_n=\hbar\omega\left(n+1/2\right)$。把传播子按
$\mathrm{e}^{-\mathrm{i}E_n T/\hbar}$ 展开，还能逐项恢复谐振子的本征波函数。
